\section{The recoverable part of a transition loss}
\label{sec:reinsurance-transition-recovery}

Group B's loss component of 20 describes the issued contracts. If the insurer
also holds reinsurance, how much of that loss is expected to be recovered?
Extend the preceding example by adding a separately recognised, qualifying
group of reinsurance contracts held. Assume it covers all contracts in B and
only B, with remaining coverage and no incurred claims, and uses general
measurement. This changes the earlier no-reinsurance assumption. The issued
obligation to policyholders remains 120, with its loss component of 20;
purchasing reinsurance does not extinguish it
\citep[paragraphs 60--64 and 75]{aasb17}.

The transition calculation has a specific timing exception. Ordinarily,
paragraph B119C permits the recovery-income treatment in paragraph 66A only
if the reinsurance contract was entered into before or when the onerous
underlying contracts were recognised. Suppose the underlying contracts'
recognition date, applying these rules, is 1 July 2020, the treaty was entered
into on 1 October 2020 and
both coverages remain in force at transition on 1 January 2022. The treaty
does not meet that ordinary condition for the original underlying recognition.
Paragraph C20A nevertheless expressly removes B119C's condition when
determining the loss-recovery component at transition. The exception does
not remove the treaty's other recognition, grouping, boundary or measurement
requirements \citep[paragraphs 61--66B, B119C and C20A]{aasb17}.

Let $q$ be the percentage of claims on the underlying contracts that the
insurer expects to recover from this reinsurance group, expressed as a
fraction. Assume $q=0.60$ for this conditional illustration. A nominal treaty
quota alone does not establish that expectation: coverage, limits, exclusions
and collectability matter. The full reinsurance measurement also includes
issuer non-performance risk and the nonfinancial risk transferred, under
paragraphs 63--64. With $K=20$ still denoting B's loss component, C20A gives
the reinsurance loss-recovery component $R$:
\[
 R=Kq=20\times 0.60=12.
\]
This 12 belongs to the reinsurance asset for remaining coverage. It is not
the issued liability, the amount of expected claims, or a valuation of the
whole reinsurance asset. If only some contracts in the underlying group were
covered, the required allocation of the losses would also have to be
established; the all-covered assumption avoids that extra calculation here
\citep[paragraphs 63--64, 66B, B119E and C20A]{aasb17}.

The complete asset still needs the treaty cash flows, discounting, transferred
risk and net cost or gain of purchasing reinsurance. The reinsurance CSM
reflects that net cost or gain; held reinsurance is not classified as onerous
under paragraph 68. Applying the issued-group CSM and loss formulas to the
held contract would therefore be wrong. Presentation likewise keeps the
issued and held-reinsurance portfolios separate
\citep[paragraphs 65--68 and 78--79]{aasb17}.

\clearpage
To extend the opening balance sheet, supply the missing total explicitly.
Assume a separate complete measurement gives a reinsurance asset of 50,
already including the loss-recovery component of 12, against an old asset of
45. These are additional conditional inputs; the multiplication above does
not derive them. Other assets remain 500 and the two issued liabilities still
increase from 220 to 230. Total assets therefore rise from $500+45=545$ to
$500+50=550$, while equity moves from $545-220=325$ to $550-230=320$.
The asset increase adds 5 to equity against the issued-liability increase
of 10, giving a net reduction of 5. An illustrative reinsurance entry debits
the new asset 50, credits the old asset 45 and credits opening equity 5.
Adding another asset of 12 would count a component of the supplied 50 twice
and overstate equity by 12. The net transition change belongs in equity
under C4(b)--(c), rather than being inferred as current recovery income
from the ordinary recognition rule \citep[paragraphs C4--C5 and C20A]{aasb17}.

The component also imposes a useful later check. If B's loss component falls
to 14 and the recoverable portion remains 60\%, paragraph B119F limits the
reinsurance loss-recovery component to $14\times0.60=8.4$. Leaving it at 12
would violate that bound. The bound alone neither selects a unique new
component nor determines the full reinsurance asset or period's profit;
the actual changes and allocation rules still matter. Similarly, $q=0$
gives no loss-recovery component and $q=1$ gives 20, without erasing the
issued obligation or determining the entire asset
\citep[paragraphs 66--67 and B119F]{aasb17}.

\clearpage
